\chapter{Introduction}
\label{chapter_introduction}

\section{Context and Motivation}

Business intelligence and analytics organize operational data so that decision
makers can describe performance, anticipate future conditions and choose an
appropriate response \cite{chen2012business}. For an inventory-based small or
medium enterprise (SME), this progression begins with familiar records: invoice
lines, product movements, purchase receipts, expenses, payments and customer
visits. These records can explain revenue and stock position, but their value is
limited when they remain divided among paper ledgers, spreadsheets and separate
applications. The owner may know that sales declined last month without knowing
which products, customers or operating conditions caused the decline, whether
the pattern is likely to continue, or which affordable action should be taken.

The decision is more demanding than producing a dashboard. Before Ramadan or
Eid, for example, a retailer must estimate demand early enough to order from a
supplier. Ordering too little creates stock-outs and forgone sales; ordering too
much immobilizes working capital and may create obsolete stock. At the same
time, a previously valuable customer may be becoming inactive, a product may
appear profitable only because costs were recorded incompletely, or a promotion
may increase volume without increasing contribution. These questions draw on
the same operational history but require different analytical tasks and a
common decision context.

Retail forecasting research identifies hierarchy, intermittency, promotions,
new products, changing assortment and the organizational use of forecasts as
central practical concerns \cite{fildes2022retail}. Evidence from large forecast
competitions also cautions against assuming that a more complex algorithm is
automatically more accurate: combinations, statistical methods and carefully
constructed baselines remain competitive across many series
\cite{makridakis2022m5}. The relevant research problem is
therefore not to attach a fashionable model to a sales file. It is to construct
a defensible path from business events to out-of-time predictions and then from
predictions to feasible, explainable actions.

\section{Bangladesh SME Context}

The intended setting is an inventory-holding retail SME in Bangladesh. This
scope is narrower than the approved thesis title, but it makes the business
entities, prediction units and decisions testable. Bangladesh-specific studies
show that ICT adoption in SMEs is shaped by perceived benefit, management and
government support, cost, complexity and the surrounding organizational and
technological conditions \cite{hoque2016ict,shahadat2023digital}. Research on
mobile financial services further shows that users value accessibility,
mobility, self-control and other affordances embedded in everyday practice
\cite{hazra2021mfs}. These findings imply that local payment channels,
low-friction data entry and owner-readable explanations are design requirements,
not optional presentation details.

This study uses a functional definition of its population rather than relying
on a regulatory size threshold that may change. A target business holds stock,
records sales and purchases, makes branch- or product-level decisions, and has
limited specialist analytics capacity. The unit of analysis may be a business,
branch, stock-keeping unit (SKU), customer or historical customer snapshot,
depending on the research question. Conclusions will be restricted to the
participating firms and the evidence actually collected.

\section{Problem Statement}

The central problem is the absence of a trustworthy and usable connection among
three activities: recording SME operations, extracting predictive insight and
selecting a feasible business action. A transaction system can preserve orders
and stock but still provide no forecast. A standalone model can report an
accuracy score but still be unusable because its features would not exist at the
decision time, its test data leaked information from training, or its output is
not linked to budget, stock, lead time and customer consent. A recommendation
can be technically valid yet operationally impossible.

The problem is decomposed into four linked gaps:

\begin{enumerate}[label=G\arabic*.]
  \item \textbf{Data foundation gap.} Sales, inventory, purchases, expenses,
  customers and payments need a consistent schema, timestamps, provenance and
  organization boundaries before they can support comparable analytics.
  \item \textbf{Prediction-validity gap.} Demand and customer-risk models must
  use historical snapshots, chronological or entity-safe splits, training-only
  transformations, strong baselines and metrics appropriate to sparse or
  imbalanced outcomes. Leakage can otherwise create persuasive but invalid
  performance \cite{kaufman2012leakage}.
  \item \textbf{Decision gap.} A forecast or risk probability does not specify
  what to order, whom to contact or whether the action is affordable. Prediction
  must be connected to costs, constraints and uncertainty; this is the
  transition from predictive to prescriptive analytics
  \cite{bertsimas2020prescriptive}.
  \item \textbf{Adoption and accountability gap.} An SME owner needs a concise
  reason, confidence indication and traceable source for a suggestion. Feature
  attribution can describe model behavior, but it neither proves causality nor
  replaces human judgment \cite{lundberg2017shap,miller2019explanation}.
\end{enumerate}

Accordingly, this thesis asks how an integrated, Bangladesh-aware framework can
transform auditable retail records into leakage-controlled forecasts and
customer insights, and then rank actions that respect the operating constraints
of an SME. The work does not presume that machine learning will defeat every
baseline. Discovering that a simple rule is more reliable or economical for a
particular task is a valid outcome.

\section{Research Aim and Objectives}

\subsection{General Objective}

The general objective is to design and evaluate an AI-powered business
analytics framework that enables Bangladeshi retail SMEs to monitor performance,
extract reliable and explainable insights, and prioritize feasible strategies.

\subsection{Specific Objectives}

The study will:

\begin{enumerate}[label=O\arabic*.]
  \item define an auditable canonical data model for sales, products, stock,
  purchases, returns, customers, suppliers, payments and expenses;
  \item construct reproducible, time-indexed preprocessing and feature pipelines
  for performance analysis, SKU demand, customer inactivity and segmentation;
  \item compare simple, statistical and machine-learning approaches under
  chronological or entity-safe validation, using both predictive and
  decision-oriented measures;
  \item quantify uncertainty and provide global and case-level supporting
  factors without presenting association as causation;
  \item formulate a constraint-aware mechanism for reorder, retention and
  operational action ranking; and
  \item evaluate technical accuracy, robustness, scalability, usability and
  potential business value on clearly separated evidence tiers, culminating in
  consenting Bangladeshi SME data.
\end{enumerate}

\section{Initial Solution Concept}

The proposed solution is \system{}: the Bangladesh SME Monitoring, Analytics,
Recommendation and Tracking framework. It is an end-to-end decision-support
process rather than a new base classifier. Business events enter a governed
operational layer; validated events are transformed into time-frozen analytical
views; candidate models generate forecasts, risks and segment profiles; and a
strategy layer removes infeasible actions before ranking the remainder. Each
displayed action carries its data cutoff, model or baseline source, uncertainty
and supporting factors. Owner response and later outcome are retained for
evaluation and monitoring.

Three analytical outputs match the approved title. \emph{Performance
evaluation} summarizes revenue, margin, stock movement and customer behavior
from reconciled records. \emph{Insight extraction} estimates future demand and
customer inactivity and identifies meaningful behavioral groups.
\emph{Strategy optimization} converts those estimates into candidates such as a
reorder quantity or retention priority, subject to cash, storage, lead time,
minimum order quantity, current stock and contact consent. Chapter
\ref{chapter_methodology} defines these components formally.

\section{Research Questions}

The thesis is guided by four questions:

\begin{description}[style=nextline,leftmargin=0pt]
  \item[RQ1.] Do calendar, price, promotion and market-context features improve
  out-of-time branch--SKU demand forecasts over seasonal and statistical
  baselines in the target SME setting?
  \item[RQ2.] How accurately and reliably can a historical customer snapshot
  identify no repeat purchase during a fixed future horizon?
  \item[RQ3.] Under the same demand and operating constraints, does a
  forecast-driven reorder policy reduce stock-out and holding cost relative to
  the participating business's current policy?
  \item[RQ4.] Do concise Bangla action explanations improve owner task
  completion, confidence and perceived usability without creating excessive
  decision time?
\end{description}

\section{Scope and Boundaries}

The research covers retail and inventory-based SMEs, structured operational
data, descriptive indicators, demand forecasting, customer inactivity,
behavioral segmentation, explanation and two strategy families: replenishment
and consent-aware customer retention. The initial geographical context is
Bangladesh; an external dataset may test method portability but cannot establish
Bangladeshi business behavior. Public price data may supply market context but
cannot substitute for invoice-level evidence.

The study excludes manufacturing scheduling, credit scoring, employee
surveillance, fully autonomous price changes and unrestricted customer
profiling. It does not claim causal effects from observational feature
attributions. It will not treat generated data as field evidence, combine
synthetic and real observations in a reported real-world score, or generalize
from one retailer to all SMEs. Implementation detail, experimental results and
final conclusions are deliberately reserved for Chapters
\ref{chapter_implementation}--\ref{chapter_conclusion}.

\section{Expected Contributions and Significance}

The expected research contribution is the evaluated integration of established
methods under constraints that matter to small retailers. Specifically, the
work is expected to provide: (i) a traceable operational-to-analytical data
design; (ii) reusable historical-snapshot protocols that prevent common forms
of leakage; (iii) a Bangladesh-aware comparison of forecasting and customer
analytics baselines; (iv) a transparent bridge from prediction to constrained
action; and (v) evidence about whether owners can correctly interpret and use
the resulting recommendations.

The practical significance is equally bounded. The proposed system is intended
to help an owner see the evidence behind a decision, compare a recommendation
with a familiar rule and retain authority over the final action. It is not
intended to remove the owner from the decision or to label a statistically
accurate output as profitable before business impact is measured.

\section{Anticipated Challenges}

The most consequential dependency is suitable primary data. The intended joint
history of invoice lines, stock, products, prices and customers is difficult to
obtain from a public Bangladesh source, and partner data may contain missing
identifiers, changing SKU codes or incomplete cost records. Technical risks
include intermittent demand, short histories, new products, stock-out censoring,
class imbalance and drift. Scientific risks include target leakage, optimistic
model selection, incompatible dataset comparisons and conclusions based on one
firm. Human risks include time pressure, varied digital skill, low trust and
over-trust in confident-looking explanations. The methodology assigns a
specific prevention or evaluation step to each of these risks.

\section{Organization of the Report}

Chapter~\ref{chapter_backStudy} develops the technical and contextual
background, reviews the verified literature and identifies the research gap.
Chapter~\ref{chapter_methodology} presents the research design, data protocol,
formal tasks, proposed framework, evaluation plan and two-semester timeline.
The Implementation and Experimental Results chapters remain named placeholders
in this initial report. The Conclusions chapter synthesizes only what can be
established from the present background study and proposed methodology; final
outcome conclusions require the later empirical evidence.
